\clearpage
\fancyhead[L]{\small IRIS.}
\SourceLeaf{480}
\begin{center}
\rule{0.55\linewidth}{0.4pt}\\[2pt]
Astronomische Nachrichten. Vol. XXVI. No. 616. pp. 241 and 255. 8 November 1847.\\[2pt]
\rule{0.55\linewidth}{0.4pt}
\end{center}
\begin{center}[\textsc{Gauss} to \textsc{Schumacher}.]\quad Göttingen, 11 October 1847.\end{center}

After a break of three full weeks, I can send you another meridian observation of \emph{Iris} today after all. The planet, however, was very difficult to observe because it was so faint.

\begin{center}
\begin{tabular}{ll}
\multicolumn{2}{l}{1847 Oct. 10.\quad $6^{\rm u}\ 42^{\rm m}\ 2^{\rm s}6$ mean time} \\
R.A. & $299^\circ\ 13'\ 32''7$ \\
Southern decl. & $14^\circ\ 19'\ 38''6$ \\
\end{tabular}
\end{center}

\textsc{Goldschmidt}’s elements still agree very well. I was able to observe only the end of the \emph{solar eclipse}.

\begin{center}
1847 Oct. 8.\quad $21^{\rm u}\ 39^{\rm m}\ 33^{\rm s}7$ Mean time.\\
Professor \textsc{Goldschmidt} $2^{\rm s}4$ later.
\end{center}

I also recorded the ingress of a spot, that is, the beginning of ingress, or first contact, at $19^{\rm u}\ 3^{\rm m}\ 40^{\rm s}8$, but the air was shimmering so strongly that it was much harder to see clearly. Professor \textsc{Goldschmidt} had $3^{\rm s}7$ less.

In a later letter the Privy Councillor added the following meridian observations of \emph{Iris}:

\begin{center}
\setlength{\tabcolsep}{4pt}
\begin{tabular}{l|r|r|r}
1847 & Mean time & Right ascension & Southern decl. \\\hline
October 11 & $6^{\rm u}\ 39^{\rm m}\ 0^{\rm s}6$ & $299^\circ\ 27'\ 3''6$ & $14^\circ\ 18'\ 39''3$ \\
\hfill 13 & $6\ 33\ 1.5$ & $299\ 55\ 18.2$ & $14\ 16\ 11.0$ \\
\hfill 14 & $6\ 30\ 4.2$ & $300\ 9\ 59.4$ & $14\ 14\ 44.5$ \\
\end{tabular}
\end{center}

\clearpage
\hypertarget{native-480-october21}{}
\begin{center}[\textsc{Gauss} to \textsc{Schumacher}.]\quad Göttingen, 21 October 1847.\end{center}

I can send you three more meridian observations of \emph{Iris}. I am also giving the observations of the 13th and 14th again here, because the new determinations of the nadir have slightly changed the declinations, and I do not know whether that change had already been applied to the numbers I sent you earlier. The figures for the 10th and 11th are unaffected.

\begin{center}
\setlength{\tabcolsep}{4pt}
\begin{tabular}{l|r|r|r}
1847 & Mean time & Right ascension & Southern decl. \\\hline
October 13 & $6^{\rm u}\ 33^{\rm m}\ 1^{\rm s}5$ & $299^\circ\ 55'\ 18''1$ & $14^\circ\ 16'\ 11''4$ \\
\hfill 14 & $6\ 30\ 4.2$ & $300\ 9\ 59.4$ & $14\ 14\ 45.1$ \\
\hfill 15 & $6\ 27\ 8.2$ & $300\ 25\ 1.5$ & $14\ 13\ 17.2$ \\
\hfill 16 & $6\ 24\ 13.9$ & $300\ 40\ 27.1$ & $14\ 11\ 50.3$ \\
\hfill 17 & $6\ 21\ 20.8$ & $300\ 56\ 12.1$ & $14\ 10\ 9.1$ \\
\end{tabular}
\end{center}

Since then it has always been overcast around the time of culmination. I regretted this all the more because it has prevented me from seeing the planet in very close conjunction with a star of the 9th magnitude from \textsc{Bessel}’s zones (20, 101 \textsc{Weisse} — on the 18th the planet must have come within almost half a minute of that star). Apart from that, the planet has always been very difficult to observe because it is so faint, but I cannot say that the closeness of the Moon on the 17th made observing any harder than it had been on the earlier evenings.